\documentclass[12pt]{ctexart}
\usepackage{geometry}
\geometry{a4paper, margin=2.5cm}

\usepackage{enumitem}
\usepackage[most]{tcolorbox}
\usepackage{xcolor}
\usepackage{hyperref}
\hypersetup{
	colorlinks=true,
	linkcolor=teal,
	urlcolor=teal
}

% ====== 颜色设置：浅绿色主题 ======
\definecolor{mygreen}{RGB}{101, 171, 139}   % 深绿（框线）
\definecolor{lightgreen}{RGB}{236, 246, 238} % 浅绿背景

% ====== 摘抄卡片样式 ======
\tcbset{
	excerptbox/.style={
		enhanced,
		sharp corners=all,
		colback=lightgreen,
		colframe=mygreen,
		coltitle=mygreen,
		boxrule=0.8pt,
		left=8pt,
		right=8pt,
		top=6pt,
		bottom=6pt,
		fonttitle=\bfseries,
		arc=3mm,
	}
}

% ====== 自动编号摘抄环境 ======
\newcounter{excerpt}
\newenvironment{excerpt}{
	\refstepcounter{excerpt}
	\begin{tcolorbox}[excerptbox, title={摘抄 \theexcerpt}]
	}{
	\end{tcolorbox}
	
	\vspace{0.4em}
	{\bfseries 感悟：}\\
}

\begin{document}
	
	% ====== 封面 ======
	\title{\Huge 读书笔记—《我身上有个不可战胜的夏天》}
	
	\author{\Large 作者：加缪/Liu}
	\date{\today}
	
	\maketitle
	\begin{center}
		\color{mygreen}
		\itshape
		\large Au milieu de l'hiver, j'apprenais enfin qu'il y avait en moi un été invincible.
	\end{center}
	\tableofcontents
	\newpage
	
	% ====== 正文 ======
	

	
	\section{反与正}
	
	\begin{excerpt}
		然而，外界能感受到的空气中所有的欢腾，以及倾泻在世间的所有喜悦，我能捕捉到的不过是映在白窗帘上婆娑的枝影。还有五束阳光，正耐心地将干草的气息注入房间。一阵微风拂过，窗帘上的影子便活了过来。当云层略过又离开太阳的瞬间，阴影中会突然跃出合金欢花瓶中耀眼的明黄。仅此而已：只需一道初生的微光，我便被一种令人眩晕的迷乱喜悦填满。
	\end{excerpt}
	这里所描述的景象也是我每天坐在书桌前所能看到的，德国这边的天气以阴天为主，偶尔洒落下来的阳光铺满桌面就让人感受到心情愉悦了。
	
	\vspace{1.5em}
	
	\begin{excerpt}
		我要说，也即将说出：重要的是活得真实，如此人性与纯真则自然显现。当我与这世界浑然一体时，何曾更真实过？未及渴望，我已满足。永恒就在这里，而我曾向往它。此刻我不再祈求幸福，只愿清醒地活着。
	\end{excerpt}
	清醒地活着就是一种幸福，跟随自己的步伐，掌控自己的节奏，去做自己喜欢的事情。
	
	\vspace{1.5em}
	
	
	\begin{excerpt}
	但若在忙碌中迷失，勤勉何尝不是一种虚度？
	\end{excerpt}
	这段话在出现在上一段话之前，若迷失则并非清醒，我们也将失去方向，长此以往我们也不知道到底在为何而忙碌与努力，如同机器一般周而复始。很多人说，做什么事情一定要有意义才去做吗？但是有无意义和迷失并无关系，无意义所带来的结果并非不好，相反可能会有其他的回报，而迷失可能会使人越陷越深。
	
	\vspace{1.5em}
	\section{总结}
	整体读完后的理解、对你的影响、未来行动。
	
\end{document}